% !PRINTED 55
\weberpage{55}
%<<CONT

so ergiebt die Ausführung der Multiplication der beiden Formeln (4), wenn man in dem Product den Coëfficienten von $y^\nu$ aufsucht:
\begin{equation}
 {D_\nu(f\varphi)\over\Pi(\nu)}
 =u_\nu v_0+u_{\nu-1}v_1+u_{\nu-2}v_2+\cdots+u_0v_\nu,
\tag{6}
\end{equation}
oder
\begin{equation}
 D_\nu(f\varphi)=\varphi D_\nu f+B_1^{(\nu)}D_{\nu-1}fD_1\varphi
 +B_2^{(\nu)}D_{\nu-2}fD_2\varphi+\cdots,
\tag{7}
\end{equation}
worin $B_1^{(\nu)},B_2^{(\nu)},\ldots$ die Binomialcoefficienten sind. In ähnlicher Weise kann man unter Anwendung der Polynomialcoefficienten die Derivirten eines Productes von mehr als zwei Factoren bilden.

Wir wollen die \gesperrt{erste Derivirte}, die wir jetzt mit $D$ statt mit $D_1$ bezeichnen, für ein Product von $n$ Factoren danach bilden. Für zwei Factoren erhalten wir nach (7):
\[
 D(f\varphi)=\varphi Df+fD\varphi=\varphi f'(x)+f\varphi'(x)
\]
und allgemein, wenn wir die Factoren mit $u_1,u_2,\ldots,u_n$, die Derivirten mit $u'_1,u'_2,\ldots,u'_n$ bezeichnen:
\begin{equation}
\begin{aligned}
D(u_1u_2\cdots u_n)
&=u'_1u_2\cdots u_n+u_1u'_2\cdots u_n+\cdots+u_1u_2\cdots u'_n,
\end{aligned}
\tag{8}
\end{equation}
oder kürzer:
\[
 {D(u_1u_2\cdots u_n)\over u_1u_2\cdots u_n}
 =\sum {D u_\nu\over u_\nu}.
\]

Wenn wir also ein Product aus linearen Factoren
\begin{equation}
 f(x)=(x-\alpha_1)(x-\alpha_2)\cdots(x-\alpha_n)
\tag{9}
\end{equation}
betrachten, so erhalten wir, da die ersten Derivirten von $x-\alpha_1$, $x-\alpha_2,\ldots,x-\alpha_n$ sämmtlich gleich $1$ sind,
\begin{equation}
\begin{aligned}
f'(x)&=(x-\alpha_2)(x-\alpha_3)\cdots(x-\alpha_n)\\
&\quad +(x-\alpha_1)(x-\alpha_3)\cdots(x-\alpha_n)+\cdots\\
&\quad +(x-\alpha_1)(x-\alpha_2)\cdots(x-\alpha_{n-1}),
\end{aligned}
\tag{10}
\end{equation}
wofür man auch setzen kann:
\begin{equation}
 f'(x)={f(x)\over x-\alpha_1}+{f(x)\over x-\alpha_2}+\cdots+{f(x)\over x-\alpha_n}.
\tag{11}
\end{equation}
Ein sehr wichtiges Resultat ergiebt sich hieraus, wenn $x$ gleich einem der Werthe $\alpha_1,\alpha_2,\ldots,\alpha_n$ gesetzt wird, nämlich
%CONT>>
